\subsection{Real radiation in invariant energy variables}
\label{sec:comp-supp-real-invariants}
The energy ordering used above is convenient for numerical checks.
For analytic integration it is useful to integrate over both photon
labels and include the identical-particle factor $1/2!$.
All invariants in this section are again measured in units of $m^2$.
In the center-of-mass frame define
\begin{equation}
 H=s-1,\qquad x=2\sqrt{s}\,\omega_1,\qquad
 y=2\sqrt{s}\,\omega_2,\qquad z=x+y-H.
\end{equation}
Momentum conservation and the final electron mass shell give
\begin{equation}
 2k_1\cdot k_2=z,\qquad
 2p'\cdot k_1=H-y,\qquad 2p'\cdot k_2=H-x.
\end{equation}
Since $0\leq1-\widehat{\bm n}_1\cdot\widehat{\bm n}_2\leq2$,
$0\leq z\leq xy/s$. Thus the exact integration region is
\begin{equation}
 0<x<H,\qquad H-x<y<\frac{s(H-x)}{s-x}.
 \label{eq:comp-real-dalitz-region}
\end{equation}
This two-dimensional region in final-state energy invariants is
called the Dalitz region. Its boundaries are rational functions. The limits $x=0$ and $y=0$
are the two soft endpoints, while $z=0$ means that the final photons
are parallel.

To obtain the measure, first integrate the three-momentum delta
function over the final electron momentum. Let
$c_{12}=\widehat{\bm n}_1\cdot\widehat{\bm n}_2$ and let $\phi$
be the azimuth of the second photon around the first. The remaining
energy delta function gives
\begin{align}
 E'&=\sqrt{1+\omega_1^2+\omega_2^2+2\omega_1\omega_2c_{12}},\\
 \dd\Phi_3&=\frac{\omega_1\omega_2}{8(2\pi)^5 E'}
 \delta(\sqrt{s}-\omega_1-\omega_2-E')
 \dd\omega_1\dd\omega_2\dd\Omega_1\dd\phi\dd c_{12}.
\end{align}
The absolute derivative of the delta-function argument with respect
to $c_{12}$ is $\omega_1\omega_2/E'$. Its integral therefore
cancels precisely the energy-dependent prefactor in this measure.
A root exists with $-1\leq c_{12}\leq1$ exactly in the Dalitz region
already derived.

For fixed $x,y$, the relative angle is now fixed. Choose a frame
attached to the two outgoing photon directions, and denote the incoming
photon direction in that frame by $\widehat{\bm n}$. The orientation
measure $\dd\Omega_1\dd\phi$ has total volume $8\pi^2$.
Rotations about the incoming direction leave the unpolarized amplitude
unchanged and integrate to $2\pi$, leaving $\dd\Omega_{\bm n}$.
Finally $\dd\omega_1\dd\omega_2=\dd x\dd y/(4s)$, so
\begin{equation}
 \dd\Phi_3=\frac{\dd x\,\dd y}{128\pi^3s}
                    \frac{\dd\Omega_{\bm n}}{4\pi}.
 \label{eq:comp-real-dalitz-measure}
\end{equation}
For example, its massless limit integrates over a triangle of area
$s^2/2$ and gives $\Phi_3=s/(256\pi^3)$.
With $\mathcal R=\overline{|\mathcal M_{\gamma\gamma}|^2}/e_{\mathrm T}^6$,
the real cross section is consequently
\begin{equation}
 \sigma_{\mathrm R}=\frac{\alpha_{\mathrm T}^3}{8sH m^2}
 \int_0^H\dd x\int_{H-x}^{s(H-x)/(s-x)}\dd y\,
 \langle\mathcal R\rangle_{\bm n},\qquad
 \langle f\rangle_{\bm n}=\int\frac{\dd\Omega_{\bm n}}{4\pi}f.
 \label{eq:comp-real-dalitz-cross-section}
\end{equation}
This formula requires the same infrared subtraction as the earlier
real contribution. Its unsubtracted integral diverges at the soft
endpoints.

\subsection{Analytic direction integrals}
Introduce three positive electron--photon invariants,
\begin{align}
 U&=2p\cdot k_1=\frac{x}{2s}
       [s+1+H\widehat{\bm n}\cdot\widehat{\bm n}_1],\\
 V&=2p\cdot k_2=\frac{y}{2s}
       [s+1+H\widehat{\bm n}\cdot\widehat{\bm n}_2],\\
 W&=2p'\cdot k=U+V-z.
 \label{eq:comp-real-angular-denominators}
\end{align}
The labels $UV$, $UW$, and $VW$ below identify pairs of these
directional denominators. In particular, $UV$ here does not mean
ultraviolet. The six electron-line orderings have denominator products
proportional to
\begin{equation}
 H(H-y),\quad H(H-x),\quad U(H-x),\quad V(H-y),
 \quad UW,\quad VW.
\end{equation}
The signs of the individual propagators follow directly from
$(p-k_i)^2-1=-2p\cdot k_i$ and
$(p'-k)^2-1=-2p'\cdot k$.
Thus all directional denominators of the spin-summed squared
amplitude are powers of $U,V,W$. Its numerator is polynomial in
these invariants, since all remaining scalar products follow from
momentum conservation.

Their affine relation gives a useful partial-fraction identity,
\begin{equation}
 \frac1{UVW}=\frac1z\left(\frac1{UW}+\frac1{VW}-\frac1{UV}\right).
 \label{eq:comp-real-three-poles}
\end{equation}
Repeated application after multiplication by further inverse powers
reduces every term with three denominator types to terms with at
most two. This operation is valid for $z>0$. Any apparent pole at
$z=0$ introduced by this reduction must be combined before taking
that limit.

The basic two-denominator direction integral has an elementary
closed form. Write $A=a_0-\bm a\cdot\widehat{\bm n}$ and
$B=b_0-\bm b\cdot\widehat{\bm n}$, where $a$ and $b$ are future
timelike four-vectors. Define
\begin{equation}
 d=a\cdot b,\qquad \Delta=d^2-a^2b^2.
\end{equation}
The bounds $A\geq a_0-|\bm a|>0$ and
$B\geq b_0-|\bm b|>0$ hold uniformly on the sphere. Their convex
combinations are bounded below by
$\min(a_0-|\bm a|,b_0-|\bm b|)>0$, so the Feynman parameter integral
can be interchanged with the direction average.
Direct polar integration of $[vA+(1-v)B]^{-2}$ gives
$[va+(1-v)b]^{-2}$, where the latter square is a Minkowski square.
Consequently,
\begin{align}
 J(a,b)&=\left\langle\frac1{AB}\right\rangle_{\bm n}
 =\int_0^1\frac{\dd v}{v^2a^2+2v(1-v)d+(1-v)^2b^2}\\
 &=\frac1{2\sqrt\Delta}\ln\frac{d+\sqrt\Delta}{d-\sqrt\Delta}.
 \label{eq:comp-real-angular-master}
\end{align}
For $\Delta>0$, the substitution $u=v/(1-v)$ evaluates the parameter
integral explicitly:
\begin{equation}
 J(a,b)=\int_0^\infty\frac{\dd u}{a^2u^2+2du+b^2}
 =\left[\frac1{2\sqrt\Delta}
 \ln\frac{a^2u+d-\sqrt\Delta}{a^2u+d+\sqrt\Delta}
 \right]_{u=0}^{u=\infty}.
\end{equation}
Future timelike vectors obey $d>0$, $\Delta\geq0$, and
$d>\sqrt\Delta$. Thus both logarithm arguments are positive, and
the primitive tends to zero at infinity. These endpoint values fix
the real logarithm and its constant. At $\Delta=0$ the quadratic
denominator is $(\sqrt{a^2}u+\sqrt{b^2})^2$, whose integral is $1/d$.
The positive bounds on $A$ and $B$ persist in a neighborhood of any
fixed future timelike pair. The integrand and its derivatives of any
finite order are then uniformly bounded on the sphere, which permits
differentiation under the integral.
Holding $\bm a,\bm b$ fixed, all higher denominator powers follow by
differentiation,
\begin{equation}
 \left\langle A^{-r}B^{-t}\right\rangle_{\bm n}
 =\frac{(-\partial_{a_0})^{r-1}(-\partial_{b_0})^{t-1}}
             {(r-1)!(t-1)!}J(a,b),\qquad r,t\geq1.
 \label{eq:comp-real-angular-powers}
\end{equation}
These derivatives are taken before substituting the physical
invariants. They are finite algebraic differentiations of the
logarithm in \eqref{eq:comp-real-angular-master}.
For a single denominator the corresponding formula is
\begin{equation}
 \langle A^{-1}\rangle_{\bm n}
 =\frac1{2|\bm a|}\ln\frac{a_0+|\bm a|}{a_0-|\bm a|},\qquad
 \langle A^{-2}\rangle_{\bm n}=\frac1{a^2}.
\end{equation}
The physical vectors corresponding to $U,V,W$ satisfy
\begin{align}
 a_U^2&=x^2/s,\qquad a_V^2=y^2/s,\qquad a_W^2=H^2/s,\\
 a_U\cdot a_V&=xy/s+H^2z/(2s),\\
 a_U\cdot a_W&=xH/s+zH(H-x)/(2s),\\
 a_V\cdot a_W&=yH/s+zH(H-y)/(2s).
 \label{eq:comp-real-angular-Gram}
\end{align}
In particular they remain timelike throughout the interior of the
Dalitz region. These identities make each scalar angular master a
closed logarithm of $s,x,y$.

The next section carries out this reduction for the full six-diagram
trace and lists the resulting coefficients. The remaining $x,y$
integrations still require the infrared subtraction.

\subsection{Explicit completion of the direction integral}
The preceding reduction can be carried out for the full six-diagram
square. We describe the algebra and give its coefficients so that
the result can be reconstructed without evaluating Dirac matrices.
The ordered photon triples are
\begin{equation}
 (a,b,c)=(0,1,2),(0,2,1),(1,0,2),(1,2,0),(2,0,1),(2,1,0).
\end{equation}
Their signed denominator products are, in this order,
\begin{equation}
 d_i=\bigl(H(H-x),H(H-y),-U(H-x),UW,-V(H-y),VW\bigr)_i.
\end{equation}
Let $T_{ij}$ be the numerator trace for the interference of
orders $i,j$, including the initial average $1/4$ and three
polarization sums $-\eta_{\mu\nu}$. Then
\begin{equation}
 \mathcal R=\sum_{i\leq j}\frac{(2-\delta_{ij})T_{ij}}{d_i d_j}.
\end{equation}
There are 21 independent terms, including all 36 interferences.
The trace reduction follows directly by moving the first contracted
gamma matrix through the intervening string. For
$\Gamma_n=\gamma^{\alpha_1}\cdots\gamma^{\alpha_n}$ in four dimensions,
\begin{equation}
 \gamma_\mu\Gamma_n\gamma^\mu
 =4(-1)^n\Gamma_n+
 2\sum_{j=1}^n(-1)^{j-1}
 \gamma^{\alpha_1}\cdots\widehat{\gamma^{\alpha_j}}\cdots
 \gamma^{\alpha_n}\gamma^{\alpha_j}.
\end{equation}
The hat denotes omission. Apply this relation to the three photon
index pairs, using cyclicity of the trace to bring each pair to
the two ends. The remaining traces contain only slashed momenta
and obey
\begin{equation}
 \operatorname{tr}(\slashed v_1\cdots\slashed v_{2n})
 =\sum_{j=2}^{2n}(-1)^j(v_1\cdot v_j)
 \operatorname{tr}(\slashed v_2\cdots
       \widehat{\slashed v_j}\cdots\slashed v_{2n}),
 \qquad \operatorname{tr}\mathbf1=4.
\end{equation}
All scalar products are fixed by the Gram matrix of $(p,k,k_1,k_2)$,
\begin{equation}
 G=\frac12\begin{pmatrix}
 2&H&U&V\\ H&0&x-U&y-V\\ U&x-U&0&z\\ V&y-V&z&0
 \end{pmatrix}.
\end{equation}
Thus this is a finite polynomial calculation with no integration
or input from a published correction factor.

After applying \eqref{eq:comp-real-three-poles} and cancelling
numerator factors, the result has 25 monomials,
\begin{equation}
 \mathcal R=\sum_{\bm e}C_{\bm e}U^{e_1}V^{e_2}W^{e_3},\qquad
 \boxed{\langle\mathcal R\rangle_{\bm n}
          =\sum_{\bm e}C_{\bm e}\mathfrak M_{\bm e}.}
 \label{eq:comp-real-angle-completed}
\end{equation}
The coefficients are listed below. The moments $\mathfrak M$ in
this formula are the following closed expressions, not unevaluated
angular integrals.

For two denominator types write $A=a_0-\bm a\cdot\widehat{\bm n}$,
$B=b_0-\bm b\cdot\widehat{\bm n}$ and abbreviate
$\alpha=a^2$, $\beta=b^2$, $d=a\cdot b$,
$\Delta=d^2-\alpha\beta$. Here $\alpha,\beta$ are auxiliary
kinematic symbols and are not the electromagnetic coupling.
Let $J$ be the logarithm in \eqref{eq:comp-real-angular-master}.
The required inverse-power moments are
\begin{align}
 \langle A^{-1}B^{-1}\rangle&=J,\\
 \langle A^{-2}B^{-1}\rangle
 &=\frac{(b_0d-a_0\beta)J+a_0d/\alpha-b_0}{\Delta},\\
 \langle A^{-1}B^{-2}\rangle
 &=\frac{(a_0d-b_0\alpha)J+b_0d/\beta-a_0}{\Delta},\\
 \langle A^{-2}B^{-2}\rangle
 &=\frac{Q_JJ+Q_0}{\Delta^2},\\
 Q_J&=2a_0b_0(\alpha\beta+2d^2)
       -d(\Delta+3\alpha b_0^2+3\beta a_0^2),\\
 Q_0&=\Delta+(2\alpha\beta+d^2)
          \left(\frac{a_0^2}{\alpha}+\frac{b_0^2}{\beta}\right)
          -6a_0b_0d.
\end{align}
For positive powers of $A$ and one denominator type $B$, define
\begin{align}
 r^2&=b_0^2-\beta,\qquad
 h=\frac{a_0b_0-d}{r^2},\qquad g=a_0-hb_0,\\
 v_\perp^2&=a_0^2-\alpha-\frac{(a_0b_0-d)^2}{r^2},\qquad
 L_B=\frac1{2r}\ln\frac{b_0+r}{b_0-r},\\
 n_0&=g^2-\frac{\beta v_\perp^2}{2r^2},\qquad
 n_1=2gh+\frac{b_0v_\perp^2}{r^2},\qquad
 n_2=h^2-\frac{v_\perp^2}{2r^2}.
\end{align}
Integration over the azimuth about $\bm b$ first gives the
polynomials $g+hB$ and $n_0+n_1B+n_2B^2$. Hence
\begin{align}
 \langle A/B\rangle&=gL_B+h,\qquad
 \langle A/B^2\rangle=g/\beta+hL_B,\\
 \langle A^2/B\rangle&=n_0L_B+n_1+n_2b_0,\\
 \langle B^{-1}\rangle&=L_B,\qquad
 \langle B^{-2}\rangle=1/\beta,\qquad
 \langle A\rangle=a_0,\qquad\langle1\rangle=1.
\end{align}
These formulas cover every monomial in
\eqref{eq:comp-real-angle-completed}. For the physical substitutions,
\begin{equation}
 (a_{U0},a_{V0},a_{W0})
 =\left(\frac{x(s+1)}{2s},\frac{y(s+1)}{2s},
                    \frac{(s-1)(2s-x-y)}{2s}\right).
\end{equation}
Their squares and scalar products are given in
\eqref{eq:comp-real-angular-Gram}.
When $\Delta$ or $r$ vanishes while the denominator vectors remain
future timelike, the expressions mean their continuous limits.
The same positive lower bounds justify taking these limits under
the direction integral. Apparent inverse powers at these points
are removable. Soft limits in which a denominator vector vanishes
instead require the infrared subtraction.

The direction integration is now complete analytically. The
remaining real contribution is a two-dimensional Dalitz integral
of explicit rational functions and logarithms, with its two soft
endpoints still to be combined with the virtual correction.
% BEGIN EXACT REAL ANGULAR COEFFICIENTS
\subsubsection{Coefficients of the integrated direction moments}
The following coefficients use $y=H-x+z$ to shorten the polynomials.
Every exponent triple not displayed has zero coefficient.
\begin{equation}
C_{-2,0,-2}=-32.
\end{equation}
\begin{equation}
C_{-2,0,-1}=\frac{8 \left(H^{2} - 2 H x - 4 H + x^{2} + 4 x - 8\right)}{- H + x}.
\end{equation}
\begin{equation}
C_{-2,0,0}=- \frac{8 \left(- H z + 4 H + x z - 4 x + 4\right)}{\left(- H + x\right)^{2}}.
\end{equation}
\begin{equation}
C_{-2,1,0}=\frac{8}{- H + x}.
\end{equation}
\begin{equation}
C_{-1,-1,0}=\frac{P_{-1,-1,0}}{z^{2} \left(- H + x\right) \left(x - z\right)}.
\end{equation}
\begin{equation}
\begin{aligned}
P_{-1,-1,0}={}& - 64 x^{2} + 32 z^{3} - 24 x z^{3} - 4 H z^{4} - 4 H^{3} z^{2} + 16 H z^{3} + 16 H^{2} z^{2}\\
& + 24 x^{2} z^{2} + 32 H z^{2} + 32 H^{2} z + 64 H x + 64 x z - 32 H x z^{2} - 8 H^{2} x z + 8 H x^{2} z.
\end{aligned}
\end{equation}
\begin{equation}
C_{-1,0,-2}=\frac{8 \left(z^{2} + 4 z - 8\right)}{z}.
\end{equation}
\begin{equation}
C_{-1,0,-1}=\frac{P_{-1,0,-1}}{H z^{2} \left(- H + x\right) \left(x - z\right)}.
\end{equation}
\begin{equation}
\begin{aligned}
P_{-1,0,-1}={}& - 64 z^{3} + 32 z^{4} - 128 H z^{3} - 64 H^{2} x - 56 H^{2} z^{3} - 32 H^{3} z - 8 x z^{4} - 4 x z^{5}\\
& - 4 H^{2} z^{4} - 4 H^{4} z^{2} - 4 x^{3} z^{3} + 4 H z^{5} + 4 H^{3} z^{3} + 8 x^{2} z^{3} + 32 x z^{3}\\
& + 32 x^{2} z^{2} + 64 H x^{2} + 128 H z^{2} + 128 H^{2} z^{2} - 64 H x z - 64 H x z^{2}\\
& - 56 H x^{2} z^{2} - 12 H^{2} x z^{3} - 12 H^{2} x^{2} z^{2} - 8 H^{2} x^{2} z + 4 H x z^{4}\\
& + 4 H x^{3} z^{2} + 8 H^{3} x z + 12 H x^{2} z^{3} + 12 H^{3} x z^{2} + 48 H^{2} x z^{2} + 64 H x z^{3}.
\end{aligned}
\end{equation}
\begin{equation}
C_{-1,0,0}=\frac{P_{-1,0,0}}{H \left(- H + x\right)^{2} \left(x - z\right)}.
\end{equation}
\begin{equation}
\begin{aligned}
P_{-1,0,0}={}& - 64 z - 16 x^{3} - 8 H^{4} - 4 x^{4} + 32 x^{2} + 64 H + 64 H^{3} + 160 H^{2} - 160 H x - 96 H^{2} x\\
& - 64 H z - 24 H^{2} z - 24 H^{2} x^{2} - 8 H^{2} z^{2} - 4 H^{3} z - 4 x^{2} z^{2} + 16 H x^{3}\\
& + 16 x^{2} z + 20 H^{3} x + 32 x z + 40 H x^{2} + 4 H x^{2} z + 12 H x z^{2} + 16 H x z.
\end{aligned}
\end{equation}
\begin{equation}
C_{-1,1,0}=- \frac{4 \left(2 H^{2} - H x + H z - 4 H + x z + 4 x - 2 z - 8\right)}{H \left(- H + x\right) \left(x - z\right)}.
\end{equation}
\begin{equation}
C_{-1,2,0}=- \frac{4 \left(- 2 H + x\right)}{H \left(- H + x\right) \left(x - z\right)}.
\end{equation}
\begin{equation}
C_{0,-2,-2}=-32.
\end{equation}
\begin{equation}
C_{0,-2,-1}=- \frac{8 \left(x^{2} - 2 x z - 4 x + z^{2} + 4 z - 8\right)}{x - z}.
\end{equation}
\begin{equation}
C_{0,-2,0}=\frac{8 \left(x z - 4 x - z^{2} + 4 z - 4\right)}{\left(x - z\right)^{2}}.
\end{equation}
\begin{equation}
C_{0,-1,-2}=\frac{8 \left(z^{2} + 4 z - 8\right)}{z}.
\end{equation}
\begin{equation}
C_{0,-1,-1}=\frac{P_{0,-1,-1}}{H z^{2} \left(- H + x\right) \left(x - z\right)}.
\end{equation}
\begin{equation}
\begin{aligned}
P_{0,-1,-1}={}& - 64 z^{3} - 8 z^{6} + 96 z^{4} - 96 H z^{3} - 96 x z^{3} - 64 H^{2} x - 56 H^{2} z^{3} - 32 H^{3} z\\
& - 16 H^{3} z^{2} - 12 x^{2} z^{4} - 8 x z^{4} + 4 x^{3} z^{3} + 8 H z^{5} + 8 x^{2} z^{3} + 16 H z^{4}\\
& + 16 x z^{5} + 32 x^{2} z^{2} + 64 H x^{2} + 96 H^{2} z^{2} + 128 H z^{2} - 64 H x z - 56 H x^{2} z^{2}\\
& - 16 H x z^{4} - 8 H^{2} x^{2} z - 4 H x^{3} z^{2} + 8 H^{3} x z + 12 H x^{2} z^{3} + 32 H x z^{3}\\
& + 80 H^{2} x z^{2}.
\end{aligned}
\end{equation}
\begin{equation}
C_{0,-1,0}=\frac{P_{0,-1,0}}{H \left(- H + x\right) \left(x - z\right)^{2}}.
\end{equation}
\begin{equation}
\begin{aligned}
P_{0,-1,0}={}& - 64 H - 64 z^{2} - 32 H^{2} - 32 x^{2} - 16 x^{3} + 4 x^{4} + 8 H^{3} + 8 z^{4} + 64 z - 96 H x\\
& - 64 H^{2} x - 40 H z^{2} - 24 x z^{3} - 16 x z^{2} - 16 x^{3} z - 8 H z^{3} - 8 H^{2} z^{2} - 4 H^{3} z\\
& + 4 H^{3} x + 8 H x^{2} + 28 x^{2} z^{2} + 32 x^{2} z + 56 H^{2} z + 96 x z + 128 H z - 4 H x^{2} z\\
& + 8 H^{2} x z + 12 H x z^{2} + 32 H x z.
\end{aligned}
\end{equation}
\begin{equation}
C_{0,0,-2}=-16.
\end{equation}
\begin{equation}
C_{0,0,-1}=\frac{P_{0,0,-1}}{H \left(- H + x\right) \left(x - z\right)}.
\end{equation}
\begin{equation}
\begin{aligned}
P_{0,0,-1}={}& - 32 z - 8 z^{2} + 4 z^{3} + 16 H^{2} + 32 H + 32 x^{2} - 32 H x - 32 x z - 12 H z^{2} + 4 x z^{2}\\
& + 8 H^{2} z + 24 H z - 4 H x z.
\end{aligned}
\end{equation}
\begin{equation}
C_{0,0,0}=\frac{P_{0,0,0}}{H^{2} \left(- H + x\right)^{2} \left(x - z\right)^{2}}.
\end{equation}
\begin{equation}
\begin{aligned}
P_{0,0,0}={}& - 32 H^{2} - 32 H^{3} - 32 z^{2} - 16 x^{4} - 72 H^{2} z^{2} - 64 H z^{2} - 48 H^{3} x - 32 H^{2} x\\
& - 32 x^{2} z - 24 x^{2} z^{2} - 20 H^{3} z^{2} - 8 H z^{3} - 4 H^{4} z - 4 H^{2} x^{3} + 4 H^{4} x\\
& + 8 x z^{3} + 32 H x^{2} + 32 H x^{3} + 32 x z^{2} + 32 x^{3} z + 32 H^{2} x^{2} + 48 H^{3} z + 64 H z\\
& + 96 H^{2} z - 120 H x^{2} z - 20 H^{2} x^{2} z - 4 H x^{2} z^{2} + 4 H x^{3} z + 20 H^{3} x z\\
& + 24 H^{2} x z^{2} + 40 H^{2} x z + 96 H x z^{2}.
\end{aligned}
\end{equation}
\begin{equation}
C_{0,1,0}=\frac{4 \left(3 H x - 3 H z + 2 H - 2 z\right)}{H \left(- H + x\right) \left(x - z\right)^{2}}.
\end{equation}
\begin{equation}
C_{1,-2,0}=- \frac{8}{x - z}.
\end{equation}
\begin{equation}
C_{1,-1,0}=\frac{4 \left(- H^{2} - H x - H z + x z + 4 x - z^{2} - 2 z + 8\right)}{H \left(- H + x\right) \left(x - z\right)}.
\end{equation}
\begin{equation}
C_{1,0,-1}=- \frac{4 \left(- H + z\right) \left(- H + 2 x - z\right)}{H \left(- H + x\right) \left(x - z\right)}.
\end{equation}
\begin{equation}
C_{1,0,0}=\frac{4 \left(- 3 H^{2} + 3 H x - 2 H + 2 z\right)}{H \left(- H + x\right)^{2} \left(x - z\right)}.
\end{equation}
\begin{equation}
C_{2,-1,0}=\frac{4 \left(H + x - z\right)}{H \left(- H + x\right) \left(x - z\right)}.
\end{equation}
% END EXACT REAL ANGULAR COEFFICIENTS

\subsection{Rationalization of the remaining logarithms}
After the completed angle average is collected, its functional
basis is particularly small,
\begin{equation}
 \langle\mathcal R\rangle=R_0+R_U L_U+R_V L_V+R_W L_W
                  +R_{UV}J_{UV}+R_{UW}J_{UW}+R_{VW}J_{VW}.
\end{equation}
Each $R$ is rational in $H,x,z$. The elementary single-denominator
functions obey $L_U=s\ln s/(xH)$ and $L_V=s\ln s/(yH)$.
The coefficients follow by inserting the moment formulas above in
\eqref{eq:comp-real-angle-completed}. The square roots in the other
four logarithms can each be rationalized before the remaining
energy integrals. This operation is performed separately in each
logarithmic term, with a common original integration domain.

For the $UV$ term, the angular Gram determinant is
\begin{equation}
 \Delta_{UV}=\frac{H^2z(4xy+H^2z)}{4s^2}.
\end{equation}
Its logarithm is a hyperbolic angle. Defining its positive
exponential $v$ fixes the substitution uniquely by $v\geq1$,
\begin{align}
 \frac{H^2z}{xy}&=\frac{(v-1)^2}{v},\qquad
 D_v=H^2v-x(v-1)^2,\\
 z&=\frac{x(H-x)(v-1)^2}{D_v},\qquad
 y=\frac{H^2v(H-x)}{D_v},\\
 \frac{\partial z}{\partial v}
 &=\frac{x(H-x)H^2(v^2-1)}{D_v^2},\qquad
 J_{UV}=\frac{2sv\ln v}{xy(v^2-1)}.
\end{align}
The exact upper Dalitz boundary gives $v=s$. Therefore
\begin{equation}
 0<x<H,\quad1<v<s,\qquad
 J_{UV}\dd z=\frac{2s\ln v}{D_v}\dd v.
 \label{eq:comp-real-UV-rational}
\end{equation}
No square root remains. At fixed $v$, the coefficient multiplying
$\ln v$ is a rational function of $x$.

For the $UW$ term the corresponding determinant and substitution
are
\begin{align}
 \Delta_{UW}&=\frac{H^2z(H-x)[4x+z(H-x)]}{4s^2},\\
 z&=\frac{x(v-1)^2}{(H-x)v},\qquad
 \frac{\partial z}{\partial v}=\frac{x(v^2-1)}{(H-x)v^2},\\
 J_{UW}&=\frac{2sv\ln v}{Hx(v^2-1)},\qquad
 J_{UW}\dd z=\frac{2s\ln v}{H(H-x)v}\dd v.
\end{align}
Its transformed domain is $1<v<s-x$. Exchanging the integration
order gives
\begin{equation}
 1<v<s,\qquad0<x<s-v.
 \label{eq:comp-real-UW-rational}
\end{equation}
Again the inner integration is rational. The $VW$ sector follows
by exchanging $x$ and $y$ before making the same transformation.

Finally the logarithm $L_W$ contains
$\sqrt{(s+1-z)^2-4s}$. Use the positive solution
\begin{equation}
 s+1-z=u+\frac{s}{u},\qquad\sqrt{s}\leq u\leq s.
\end{equation}
The original Dalitz bounds at fixed $z$ become
\begin{equation}
 s-u\leq x\leq s-\frac{s}{u},\qquad
 L_W|\dd z|=\frac{s}{Hu}\ln\frac{u^2}{s}\dd u.
 \label{eq:comp-real-W-rational}
\end{equation}
Thus every term has a rational inner integral with algebraic
boundaries, followed by one logarithmic integration. The square
roots and their signs have been fixed by positive physical energies
throughout.

These separated logarithmic terms must not be assigned unrelated
endpoint cutoffs. For example, a common center-of-mass soft boundary
$\omega_i>\varepsilon$ means $x,y>\delta$, where
$\delta=2\sqrt{s}\,\varepsilon$. In
\eqref{eq:comp-real-UV-rational} it becomes
\begin{equation}
 \delta<x<\frac{H^2v(H-\delta)}{H^2v-\delta(v-1)^2}.
\end{equation}
In \eqref{eq:comp-real-UW-rational} the same boundary is
$x>\delta$ together with
$H-x+x(v-1)^2/[(H-x)v]>\delta$.
The finite terms generated by these transformed boundaries belong
to the cross section. Replacing them by arbitrary small limits
would change the finite answer. Apparent poles at $z=0$ introduced
by the partial fractions must likewise cancel after the rational
and logarithmic terms are combined. The transformations must therefore be combined with the local
Gram subtraction and the common regulator matching developed below.

\subsection{Removal of artificial Gram singularities}
The original angular average is regular when $z=0$ at fixed
$0<x<H$. Individual logarithmic coefficients contain $z^{-2}$
because of the partial-fraction reduction. To remove these
artificial singularities locally, write $a=H-x$ and define
\begin{align}
 T_{UV}(z)&=\frac{s}{xa}-\frac{s(6x+H^2)}{6x^2a^2}\,z,\\
 T_{UW}(z)&=\frac{s}{xH}-\frac{sa}{6Hx^2}\,z,\\
 T_{VW}(z)&=\frac{s}{aH}-\frac{s(x+6)}{6Ha^2}\,z.
\end{align}
These are the constant and linear Taylor terms of the corresponding
$J$ functions, including $y=a+z$. They follow from
\begin{equation}
 \frac{\operatorname{arcosh}(1+q)}{\sqrt{(1+q)^2-1}}
       =1-\frac q3+\mathcal O(q^2).
\end{equation}
Replace each $R_{ij}J_{ij}$ by
$R_{ij}(J_{ij}-T_{ij})$, and add the subtracted pieces to
\begin{equation}
 R_0^{\mathrm{reg}}=R_0+
             R_{UV}T_{UV}+R_{UW}T_{UW}+R_{VW}T_{VW}.
\end{equation}
This is an exact algebraic identity. Each difference $J-T$ is
$\mathcal O(z^2)$. On substituting the explicit coefficients,
the reduced rational denominator of $R_0^{\mathrm{reg}}$ has no
factor $z$. Thus all these contributions can be integrated across
$z=0$ without assigning separate unphysical endpoint cutoffs.
This operation does not remove either physical soft endpoint.

\subsection{Analytic regulation of the soft endpoints}
A Bose-symmetric analytic regulator makes the remaining endpoint
operations independent of frame-dependent energy boundaries.
Multiply the real phase-space integrand by $(xy)^\eta$, initially
with positive real $\eta$ large enough for convergence, and then
continue meromorphically to $\eta=0$. Since $x,y>0$ in the physical
region, the real logarithm defines this factor without a branch
ambiguity. This is an intermediate integration device. The virtual
amplitude and its infrared matching still use the photon mass
$\lambda$.

Define the dimensionless regulated real coefficient by
\begin{equation}
 \mathcal C_{\mathrm R}(\eta)=\frac1{8sH}
 \int_0^H\dd x\int_{H-x}^{s(H-x)/(s-x)}\dd y\,
     (xy)^\eta\langle\mathcal R\rangle
 =\frac{A(s)}{\eta}+F(s)+\mathcal O(\eta).
 \label{eq:comp-real-analytic-regulator}
\end{equation}
The regulator is applied to the complete trace, including the
identical-photon factor already present in its normalization.
At one soft endpoint $x=2\sqrt{s}\,\omega$ while $y\to H$.
Consequently $(xy)^\eta\to(2\sqrt{s}H\omega)^\eta$.
The other endpoint gives the same contribution. Each carries
one half of the usual eikonal result because of $1/2!$, so together
\begin{equation}
 A(s)=\frac1s\int_{-1}^1\dd c\,
             \mathcal B(s,c)[b_e(s,c)-1].
 \label{eq:comp-real-pole-residue}
\end{equation}
The functions $\mathcal B$ and $b_e$ are the Born square and the
relative-velocity factor defined earlier.

This relation also fixes the finite matching constant, rather than
only the pole. Put $C=2\sqrt{s}H$. Integration of the leading soft
term below a center-of-mass energy $\varepsilon$ gives
\begin{equation}
 A(s)C^\eta\int_0^\varepsilon\omega^{\eta-1}\dd\omega
 =\frac{A(s)}\eta+A(s)\ln(C\varepsilon)+\mathcal O(\eta).
\end{equation}
Thus the real coefficient above that energy boundary is
$F(s)-A(s)\ln(C\varepsilon)+o(1)$.
The photon-mass soft term derived previously is
\begin{equation}
 \mathcal C_{\mathrm{soft}}(\varepsilon,\lambda)
 =A(s)\ln\frac{2\varepsilon}{\lambda}
                  -\frac1s\int_{-1}^1\mathcal B K\,\dd c.
\end{equation}
Adding it to the virtual contribution
$\mathcal C_{\mathrm{virt}}=(4s)^{-1}\int\operatorname{Re}\mathcal T\,\dd c$
therefore gives
\begin{equation}
 \mathcal C_{\mathrm{total}}
 =\mathcal C_{\mathrm{virt}}+F(s)
   +A(s)\ln\frac1{\sqrt{s}H\lambda}
   -\frac1s\int_{-1}^1\mathcal B K\,\dd c.
 \label{eq:comp-real-regulator-matching}
\end{equation}
The coefficient of $\ln\lambda$ cancels the virtual coefficient
$A(s)\ln\lambda$. Both the analytic pole and the temporary energy
boundary have disappeared. Equation~\eqref{eq:comp-real-regulator-matching}
fixes how the finite part $F(s)$ must be combined with the earlier
photon-mass calculation. The following analytic primitives and endpoint formulas evaluate
this finite part without phase-space quadrature.

\subsection{A complete logarithmic-sector primitive}
The part proportional to $L_W$ can be integrated further without
any phase-space quadrature. Its rational coefficient has the simple
partial-fraction form
\begin{equation}
 R_W=C_0(z)+C_1(z)\left(\frac1{x-z}-\frac1{x-H}\right),
\end{equation}
where, with $D_z=(H-z)^2-4z$,
\begin{align}
 C_0(z)&=-\frac8{HD_z}
       [H^2z-4H^2-2Hz^2+12Hz+z^3-8z^2+16z],\\
 C_1(z)&=\frac4{HD_z}\bigl[
       -H^3z-4H^3+3H^2z^2+8H^2z-8H^2\\
 &\hspace{25mm}{}-3Hz^3+2Hz^2+32Hz+z^4-6z^3+32z\bigr].
\end{align}
At fixed $z$, the exact limits are
$x_\pm=(H+z\pm\sqrt{D_z})/2$. Their substitution into the
rational primitive gives
\begin{equation}
 \int_{x_-}^{x_+}R_W\dd x
 =C_0(z)\sqrt{D_z}+2C_1(z)
       \ln\frac{H-z+\sqrt{D_z}}{H-z-\sqrt{D_z}}.
\end{equation}
Use $z=s+1-u-s/u$ as above. Define two rational functions
\begin{equation}
 A(u)=\frac{s}{Hu}C_0(s+1-u-s/u)(u-s/u),\qquad
 B(u)=\frac{2s}{Hu}C_1(s+1-u-s/u).
\end{equation}
The remaining integrand of this sector is exactly
\begin{equation}
 A(u)\ln\frac{u^2}{s}
 +B(u)\ln\frac{u^2}{s}\ln\frac{u(u-1)}{s-u},\qquad
 \sqrt{s}<u<s.
 \label{eq:comp-real-W-last-integrand}
\end{equation}
Only simple real logarithms occur. The real branches are fixed by
$u>1$ and $s-u>0$.

For a compact analytic primitive put $r=\sqrt{s}>1$. Denote by
$G_r(a_1,\ldots,a_n;u)$ the generalized polylogarithm based at $r$,
with derivative
\begin{equation}
 \frac{\dd}{\dd u}G_r(a_1,\ldots,a_n;u)
 =\frac{G_r(a_2,\ldots,a_n;u)}{u-a_1},\qquad G_r(;u)=1.
\end{equation}
The integration path is the real interval from $r$ to $u$. At a
basepoint letter $r$ use the continuous or combined tangential
limit. All other letters below are $0,1,-r,s$, which lie outside
the open interval. Thus the functions that occur are real there.
In particular, suppressing $r,u$ in the notation,
\begin{equation}
 \ln(u^2/s)=2G(0),\qquad
 \ln[u(u-1)/(s-u)]=G(0)+G(1)-G(s).
\end{equation}
Shuffle multiplication therefore reduces
\eqref{eq:comp-real-W-last-integrand} to a finite weight-two
integrand. Its primitive has weight at most three and is
\begin{align}
 \mathcal F_W={}&-\frac{16(r+1)^2}{r-1}
       [G(-r,0)+G(-r,1)-G(-r,s)]\\
 &+\frac{16(r-1)^2}{r+1}
       [G(r,0)+G(r,1)-G(r,s)]\\
 &-\frac{32r^2(2r^2+1)}{(r^2-1)^2}
       [2G(0,0,0)+G(0,0,1)+G(0,1,0)-G(0,0,s)-G(0,s,0)]\\
 &+A_{00}G(0,0)+A_{01}[G(0,1)-G(0,s)]
        +A_{10}G(1,0)+A_{s0}G(s,0)\\
 &+A_0G(0)+A_1[G(1)-G(s)]+A_{\mathrm c}.
 \label{eq:comp-real-W-GPL-primitive}
\end{align}
The rational coefficients below are obtained by partial fractions
and integration by parts. They contain no fitted constants.
Differentiating this finite expression yields
\eqref{eq:comp-real-W-last-integrand}.
The upper soft endpoint must still be combined with the other
sectors in the common analytic regulator. This primitive alone
is not the finite inclusive real correction.
% BEGIN REAL W GPL COEFFICIENTS
\begin{equation}
A_{00}=\frac{N_{00}}{u^{2} \left(r - 1\right)^{2} \left(r + 1\right)^{2} \left(r - u\right) \left(r + u\right)}.
\end{equation}
\begin{equation}
\begin{aligned}
N_{00}={}& - 16 r^{8} + 64 u^{4} - 192 r^{2} u^{3} - 112 r^{4} u^{4} - 48 r^{6} u^{2} - 16 r^{2} u^{6}\\
& - 8 r^{6} u^{4} + 8 r^{8} u^{2} + 32 r^{6} u + 32 r^{8} u + 32 r^{2} u^{5} + 32 r^{4} u^{5}\\
& + 56 r^{4} u^{2} + 64 r^{4} u^{3} + 72 r^{2} u^{4}.
\end{aligned}
\end{equation}
\begin{equation}
A_{01}=\frac{N_{01}}{u^{2} \left(r - 1\right)^{2} \left(r + 1\right)^{2} \left(r - u\right) \left(r + u\right)}.
\end{equation}
\begin{equation}
\begin{aligned}
N_{01}={}& - 8 r^{8} + 32 u^{4} - 96 r^{2} u^{3} - 88 r^{4} u^{4} - 8 r^{2} u^{6} + 8 r^{6} u^{2} + 16 r^{6} u\\
& + 16 r^{8} u + 16 r^{2} u^{5} + 16 r^{4} u^{5} + 32 r^{4} u^{3} + 64 r^{2} u^{4}.
\end{aligned}
\end{equation}
\begin{equation}
A_{10}=\frac{N_{10}}{u^{2} \left(r - 1\right)^{2} \left(r + 1\right)^{2} \left(r - u\right) \left(r + u\right)}.
\end{equation}
\begin{equation}
\begin{aligned}
N_{10}={}& - 8 r^{8} - 96 r^{2} u^{3} - 56 r^{4} u^{4} - 24 r^{6} u^{2} - 8 r^{2} u^{6} - 8 r^{8} u^{2}\\
& + 8 r^{4} u^{2} + 8 r^{6} u^{4} + 16 r^{6} u + 16 r^{8} u + 16 r^{2} u^{5} + 16 r^{4} u^{5}\\
& + 32 r^{2} u^{2} + 32 r^{4} u^{3} + 56 r^{2} u^{4}.
\end{aligned}
\end{equation}
\begin{equation}
A_{s0}=\frac{N_{s0}}{u^{2} \left(r - 1\right)^{2} \left(r + 1\right)^{2} \left(r - u\right) \left(r + u\right)}.
\end{equation}
\begin{equation}
\begin{aligned}
N_{s0}={}& - 32 u^{4} + 8 r^{8} - 56 r^{4} u^{2} - 32 r^{4} u^{3} - 16 r^{6} u - 16 r^{8} u - 16 r^{2} u^{5}\\
& - 16 r^{4} u^{5} - 8 r^{2} u^{4} - 8 r^{8} u^{2} + 8 r^{2} u^{6} + 8 r^{6} u^{4} + 24 r^{4} u^{4}\\
& + 56 r^{6} u^{2} + 96 r^{2} u^{3}.
\end{aligned}
\end{equation}
\begin{equation}
A_{0}=\frac{12 r^{2} \left(r^{4} u + 2 r^{4} + 2 r^{2} u^{2} + 6 r^{2} + 6 u^{2} - u\right)}{u \left(r - 1\right)^{2} \left(r + 1\right)^{2}}.
\end{equation}
\begin{equation}
A_{1}=\frac{N_{1}}{u^{2} \left(r - 1\right)^{2} \left(r + 1\right)^{2}}.
\end{equation}
\begin{equation}
\begin{aligned}
N_{1}={}& - 4 r^{6} - 32 r^{4} u^{2} - 12 r^{2} u^{2} - 12 r^{6} u^{2} - 4 r^{2} u^{4} + 16 r^{4} u + 16 r^{6} u\\
& + 16 r^{2} u^{3} + 16 r^{4} u^{3}.
\end{aligned}
\end{equation}
\begin{equation}
A_{\mathrm c}=\frac{4 r^{2} \left(r - u\right) \left(r + u\right) \left(7 r^{2} + 17\right)}{u \left(r - 1\right)^{2} \left(r + 1\right)^{2}}.
\end{equation}
% END REAL W GPL COEFFICIENTS

\subsection{The sector proportional to the energy logarithm}
The single denominators $U,V$ give another sector that can be
integrated completely in the interior. Substituting
$L_U=s\ln s/(xH)$ and $L_V=s\ln s/(yH)$ first combines their
coefficients into a rational function. At fixed $z$ its poles in
$x$ are only $0,z,H,H+z$, with orders at most three. The exact
limits $x_\pm$ therefore produce rational functions and the two
logarithms
\begin{equation}
 \ln\frac{x_+}{x_-}
   =\ln\frac{s(u-1)}{u(s-u)},\qquad
 \ln\frac{x_+-z}{x_--z}
   =\ln\frac{u(u-1)}{s-u}.
\end{equation}
After $z=s+1-u-s/u$, no square root remains. The final primitive
can be written with the same basepoint $r=\sqrt{s}$ as above.
Define the conjugate letters
\begin{equation}
 a_\pm=1\pm\ii\sqrt H,\qquad H=s-1>0.
\end{equation}
They are the roots of $u^2-2u+s$. Both lie off the real path,
and complex conjugation exchanges their GPLs. Thus their sum
below is manifestly real, with no choice based on numerical values.
The primitive of this entire logarithmic sector is $\ln s$ times
\begin{align}
 \mathcal F_{\ln s}={}&\frac{96s^2}{H^2}G(0,0)
 +\frac{32s(H^2-4H-6)}{H^2}[G(1,0)+G(s,0)]\\
 &+\frac{32s(H^2-H-3)}{H^2}\biggl[
 G(0,1)-G(0,s)\\
 &\hspace{18mm}{}-\sum_{a=a_\pm}
       \{G(a,0)+G(a,1)-G(a,s)\}\biggr]\\
 &+P_0G(0)+P_1[G(1)-G(s)]+P_{\mathrm c}.
 \label{eq:comp-real-log-s-primitive}
\end{align}
There are fourteen GPL words, all of weight at most two.
The rational coefficients are displayed below. Exact symbolic
differentiation of every GPL coefficient recovers the rational
and logarithmic integrand after the $x$ integration. This is an
independent check of the primitive, rather than a fit to a total
cross section. As in the previous sector, its soft endpoint is to
be combined in the common regulator before the finite part is taken.
% BEGIN REAL LOGS GPL COEFFICIENTS
\begin{equation}
P_{0}=\frac{Q_{0}}{H^{2} u^{2} \left(H + 1\right)}.
\end{equation}
\begin{equation}
\begin{aligned}
Q_{0}={}&12 - 48 u - 48 u^{3} - 32 u^{2} + 12 H^{4} + 12 u^{4} + 48 H + 48 H^{3} + 72 H^{2}\\
& - 136 H u - 120 H^{2} u - 88 H u^{3} - 56 H u^{2} - 32 H^{2} u^{3} - 24 H^{3} u\\
& + 8 H^{4} u + 8 H^{3} u^{3} + 12 H^{2} u^{4} + 24 H u^{4} + 28 H^{2} u^{2}\\
& + 36 H^{4} u^{2} + 80 H^{3} u^{2}.
\end{aligned}
\end{equation}
\begin{equation}
P_{1}=\frac{Q_{1}}{H^{4} u^{2} \left(H + 1\right)}.
\end{equation}
\begin{equation}
\begin{aligned}
Q_{1}={}& -112 - 1696 u^{2} - 796 H^{2} - 624 H^{3} - 480 H - 216 H^{4} - 112 u^{4} - 16 H^{5}\\
& + 4 H^{6} + 960 u + 960 u^{3} - 7448 H^{2} u^{2} - 5856 H u^{2} - 4152 H^{3} u^{2}\\
& - 932 H^{4} u^{2} - 256 H u^{4} - 172 H^{2} u^{4} - 112 H^{5} u^{2} - 36 H^{6} u^{2}\\
& - 24 H^{3} u^{4} + 4 H^{4} u^{4} + 24 H^{6} u + 24 H^{5} u^{3} + 224 H^{4} u^{3}\\
& + 248 H^{5} u + 1272 H^{3} u^{3} + 1496 H^{4} u + 2816 H u^{3} + 2928 H^{2} u^{3}\\
& + 3776 H u + 4200 H^{3} u + 5744 H^{2} u.
\end{aligned}
\end{equation}
\begin{equation}
P_{\mathrm c}=\frac{Q_{\mathrm c}}{H^{3} u^{2} \left(H + 1\right)}.
\end{equation}
\begin{equation}
\begin{aligned}
Q_{\mathrm c}={}& -56 - 368 u^{3} - 272 H^{2} - 206 H - 144 H^{3} - 16 H^{4} + 6 H^{5} + 56 u^{4} + 368 u\\
& - 880 H u^{3} - 660 H^{2} u^{3} - 136 H^{3} u^{3} - 6 H^{3} u^{4} - 4 H^{5} u\\
& + 4 H^{4} u^{3} + 28 H^{2} u^{4} + 94 H u^{4} + 132 H^{4} u + 796 H^{3} u + 1248 H u\\
& + 1540 H^{2} u.
\end{aligned}
\end{equation}
% END REAL LOGS GPL COEFFICIENTS



\subsection{Factorization of both soft endpoints in the pair sector}
The $UW$ sector admits a rectangular domain in which the analytic
regulator factorizes at both soft endpoints. Starting from
\eqref{eq:comp-real-UW-rational}, put
\begin{align}
 a&=H-x,\qquad v=1+at,\qquad0<a<H,\quad0<t<1,\\
 z&=\frac{a(H-a)t^2}{1+at},\qquad
 y=\frac{a[1+Ht^2+at(1-t)]}{1+at},\qquad
 \left|\frac{\partial(x,v)}{\partial(a,t)}\right|=a.
\end{align}
The product that regulates the two soft endpoints is therefore
\begin{equation}
 xy=a(H-a)f(a,t),\qquad
 f(a,t)=\frac{1+Ht^2+at(1-t)}{1+at}>0.
 \label{eq:comp-real-UW-reg-factor}
\end{equation}
Both endpoints are now ordinary powers of a single variable.
The regulator introduces no new physical branch.

Define two rational functions by the previously specified
coefficient $R_{UW}$ and Taylor subtraction $T_{UW}$,
\begin{align}
 L(a,t)&=a\,R_{UW}(x,z)\frac{2s}{H(H-x)v}
                  \bigg|_{x=H-a,\,v=1+at},\\
 T(a,t)&=a\,R_{UW}(x,z)T_{UW}(z)
                   \frac{\partial z}{\partial v}
                  \bigg|_{x=H-a,\,v=1+at},\\
 F(a,t)&=L(a,t)\ln(1+at)-T(a,t).
\end{align}
The subtraction sign in $F$ is fixed by the local removal of the
artificial Gram singularity. The complementary rational term is
already included in $R_0^{\mathrm{reg}}$.
The only poles of $L,T$ in the complex $a$ plane are
\begin{equation}
 a=0,\quad H,\quad-1/t,\quad-2/t,\quad-1/[t(1-t)].
\end{equation}
For $0<t<1$, the last three are negative and lie outside the
integration interval. Their orders are finite, at most five.
There is no polynomial denominator of degree higher than one
after partial fractions in $a$.

For clarity, the endpoint conversion can be written without a
hypergeometric expansion. Suppose
\begin{equation}
 F(a,t)=\frac{f_1(t)}a+\mathcal O(1),\qquad
 F(H-b,t)=\frac{g_2(t)}{b^2}+\frac{g_1(t)}b+\mathcal O(1).
\end{equation}
For the explicit coefficients above, $f_1(t)=0$ identically.
The Taylor subtraction has therefore also made the $a=0$ soft
corner finite in this particular sector. The other endpoint has
the stated maximum pole orders.
Let $\operatorname{FP}_a\int_0^H F\dd a$ denote the constant term
of the integral from $\delta$ to $H-\delta$ after all powers and
logarithms of $\delta$ are subtracted. Expansion of
\eqref{eq:comp-real-UW-reg-factor} at each endpoint gives
\begin{align}
 \int_0^H[xy]^\eta F(a,t)\dd a
 ={}&\frac{f_1(t)+g_1(t)}\eta
       +\operatorname{FP}_a\int_0^H F(a,t)\dd a\\
 &+f_1(t)\ln[H(1+Ht^2)]+g_1(t)\ln H\\
 &+g_2(t)\left[-\frac1H+\frac{t^2}{1+Ht}\right]
       +\mathcal O(\eta).
 \label{eq:comp-real-UW-eta-finite}
\end{align}
For example, the last term follows because the term linear in
$b$ in the logarithm of the regulator's smooth factor is
$-1/H+t^2/(1+Ht)$. Its multiplication by $\eta$ and a double
endpoint pole produces a finite term. Dropping that term would
change the answer even though the infrared pole remained correct.

Every primitive required on the right-hand side is elementary or
a real dilogarithm. To see this explicitly, define
\begin{equation}
 E_n(a;c)=\begin{cases}
 \ln|a-c|,&n=1,\\
 -(a-c)^{1-n}/(n-1),&n\geq2.
 \end{cases}
\end{equation}
For a simple logarithmic pole, let $w=1+at$ and $d=1+ct$.
When $c=-2/t$ or $c=-1/[t(1-t)]$, one has $d<0$, and a primitive is
\begin{equation}
 I_1(a;c)=\ln w\ln(1-w/d)+\operatorname{Li}_2(w/d).
\end{equation}
Both logarithms and the dilogarithm are real. For $c=H$ the same
formula applies with $0<w/d<1$ in the interior. The exceptional
values are
\begin{equation}
 I_1(a;0)=-\operatorname{Li}_2(-at),\qquad
 I_1(a;-1/t)=\tfrac12\ln^2(1+at).
\end{equation}
In each case $\partial_a I_1=\ln(1+at)/(a-c)$.
Higher poles follow by one integration by parts. With $b=-1/t$
and $n\geq2$,
\begin{align}
 I_n(a;c)&=-\frac{\ln(1+at)}{(n-1)(a-c)^{n-1}}
                  +\frac{Q_{n-1}(a;c,b)}{n-1},\\
 Q_m(a;c,b)&=\frac{E_1(a;b)}{(b-c)^m}
            -\sum_{j=1}^m\frac{E_j(a;c)}{(b-c)^{m-j+1}}
                         \quad(c\ne b),\\
 Q_m(a;b,b)&=E_{m+1}(a;b).
\end{align}
Thus partial fractions of the explicitly defined rational functions
$L,T$ give the finite part in
\eqref{eq:comp-real-UW-eta-finite} as a finite sum of boundary values
of $I_n$ and $E_n$. No numerical integration enters this inner
energy integration.

The nontrivial negative dilogarithm argument at $a=H$ is
$-(1-t)(1+Ht)/t$. Its complementary argument contains
\begin{equation}
 1+Ht-Ht^2,
 \qquad t_\pm=\frac12\left(1\pm\sqrt{1+4/H}\right).
\end{equation}
Here $t_-<0$ and $t_+>1$. The square root $\sqrt{H(H+4)}$
therefore follows directly from the remaining integration alphabet.
It has not been inferred from a published expression.
The inner integration leaves one angular variable $t$. We first
write its finite integrand explicitly, including its regulated
endpoint terms, so that the remaining primitive can be constructed
from its actual poles and logarithms.

\paragraph{Closed inner finite part.} The following boundary formulas
make the inner result a finite expression of logarithms and
dilogarithms, with no integration operation left. Put
$\ell=\ln(1+Ht)$ and write the partial fractions
\begin{equation}
 L=l_0+l_1a+\sum_{c,n}\frac{L_{cn}}{(a-c)^n},\qquad
 T=\sum_{j=0}^3t_j a^j+\sum_{c,n}\frac{T_{cn}}{(a-c)^n}.
\end{equation}
The pole set is the five-element set displayed above.
For a pole of order $m_c$, its coefficient is fixed algebraically,
\begin{equation}
 L_{cn}=\left.\frac1{(m_c-n)!}
  \left(\frac{\partial}{\partial a}\right)^{m_c-n}
       [(a-c)^{m_c}L(a,t)]\right|_{a=c},
\end{equation}
with the same formula for $T_{cn}$. The polynomial parts follow by
ordinary polynomial division. Define $\mathcal I_n(c)$ and
$\mathcal E_n(c)$ to be the finite boundary differences of the
explicit primitives $I_n,E_n$ above. At a negative pole they are
ordinary differences between $a=H$ and $a=0$. At the two physical
endpoints the required values are
\begin{align}
 \mathcal E_1(0)&=\ln H,\qquad
 \mathcal E_1(H)=-\ln H,\qquad
 \mathcal E_2(0)=\mathcal E_2(H)=-H^{-1},\\
 \mathcal I_1(0)&=-\operatorname{Li}_2(-Ht),\\
 \mathcal I_2(0)&=-\ell/H+t\ln H-t\ell+t,\\
 \mathcal I_1(H)&=-\ell\ln H-\operatorname{Li}_2(-Ht)-\ell^2/2,\\
 \mathcal I_2(H)&=-\frac{t}{1+Ht}(\ln H+\ell+1).
\end{align}
The two polynomial logarithmic moments are
\begin{align}
 P_0&=\frac{(1+Ht)\ell-Ht}{t},\\
 P_1&=\frac{[(1+Ht)^2/2-(1+Ht)]\ell
              -(1+Ht)^2/4+(1+Ht)-3/4}{t^2}.
\end{align}
The exact inner coefficient of $\eta^0$ is consequently
\begin{align}
 F_{UW}(t)={}&l_0P_0+l_1P_1
       -\sum_{j=0}^3\frac{t_jH^{j+1}}{j+1}
       +\sum_{c,n}[L_{cn}\mathcal I_n(c)-T_{cn}\mathcal E_n(c)]\\
 &+g_1(t)\ln H+g_2(t)\left[-H^{-1}+\frac{t^2}{1+Ht}\right].
 \label{eq:comp-real-UW-inner-closed}
\end{align}
This formula uses only finite algebraic differentiations and the
real functions already displayed. The corresponding $\eta^{-1}$
coefficient is $g_1(t)$.

For the last integration, every dilogarithm in this result has a
fixed-letter GPL representation. Put $b=-1/H$, $d=-2/H$, and
let $t_\pm$ be the two real roots above. With GPLs now based at zero,
\begin{align}
 \operatorname{Li}_2(-Ht)&=-G(0,b;t),\\
 \operatorname{Li}_2(-1-Ht)+\frac{\pi^2}{12}
       &=-\ln2\,G(b;t)-G(b,d;t).
\end{align}
\begin{equation}
\begin{aligned}
 &\operatorname{Li}_2\left(-\frac{(1-t)(1+Ht)}t\right)
 -\operatorname{Li}_2\left(-\frac{1-t}t\right)\\
 &\qquad=G(b,0;t)
          -\sum_{c=t_\pm}[G(1,c;t)+G(b,c;t)-G(0,c;t)].
\end{aligned}
\end{equation}
The last identity follows by differentiating both sides and using
their common zero limit as $t\to0^+$. It fixes its integration
constant as well as its branches. This gives a weight-two GPL
integrand with rational coefficients for the final energy integral.

\subsection{Closed construction for the regular rational sector}
The rational term $R_0^{\mathrm{reg}}$ also has an analytic primitive
after both Dalitz integrations. The compact construction below
specifies its coefficients by finite algebraic operations, without
introducing a new integral. At fixed $z$, its partial fractions are
\begin{equation}
 R_0^{\mathrm{reg}}=p_0(z)+
          \sum_{\rho\in\mathcal P}\sum_{n=1}^{m_\rho}
               \frac{c_{\rho n}(z)}{(x-\rho)^n},
\end{equation}
where there are sixteen nonzero pole terms and
\begin{equation}
 \mathcal P=\left\{0,z,H,H+z,
              \frac{Hz}{z-4},\ H+z-\frac{Hz}{z-4}\right\}.
\end{equation}
Every order $m_\rho$ is at most three. A uniform formula, including
vanishing coefficients when the actual order is smaller, is
\begin{align}
 c_{\rho n}(z)&=\left.\frac{1}{(3-n)!}
   \frac{\partial^{3-n}}{\partial x^{3-n}}
       [(x-\rho)^3R_0^{\mathrm{reg}}(x,z)]\right|_{x=\rho},
       \quad n=1,2,3,\\
 p_0(z)&=\lim_{x\to\infty}R_0^{\mathrm{reg}}(x,z).
\end{align}
The limit exists because its polynomial part is constant.
This supplies an explicit algebraic definition from the twenty-five
trace coefficients and the displayed Taylor subtractions.

Let $z=s+1-u-s/u$, $x_-=s-u$, $x_+=s-s/u$, and
$j=1-s/u^2$. Define four rational functions by
\begin{align}
 R(u)&=j\left[p_0(x_+-x_-)
  +\sum_{\rho}\sum_{n=2}^{m_\rho}
      \frac{c_{\rho n}}{1-n}
       \{(x_+-\rho)^{1-n}-(x_--\rho)^{1-n}\}\right],\\
 K_1(u)&=j[-c_{0,1}+c_{H+z,1}],\\
 K_2(u)&=j[-c_{H,1}+c_{0,1}+c_{z,1}-c_{H+z,1}],\\
 K_3(u)&=j[c_{Hz/(z-4),1}-c_{H+z-Hz/(z-4),1}].
\end{align}
All quantities on the right are evaluated at $z=s+1-u-s/u$.
The exact $x$ boundary differences then give the final integrand
\begin{equation}
 R(u)+K_1(u)L_1(u)+K_2(u)L_2(u)+K_3(u)L_3(u),
\end{equation}
where
\begin{align}
 L_1&=\ln(u^2/s),\qquad
 L_2=\ln[u(u-1)/(s-u)],\\
 L_3&=\ln\frac{(u-1)(u+s)^2}{u(u+1)^2(s-u)}.
\end{align}
Each logarithm vanishes at $u=\sqrt{s}$. In terms of GPLs based
there, the additional logarithm is
\begin{equation}
 L_3=G(1)+2G(-s)-G(0)-2G(-1)-G(s).
\end{equation}

For precision, define the finite algebraic operation
$\mathcal P[R;w]$ on a rational function $R$ and an empty or
one-letter word $w$. Decompose $R$ into polynomial and pole terms.
On a simple pole its value is
\begin{equation}
 \mathcal P[c/(u-a);w]=cG(a,w;u).
\end{equation}
For a pole of order $n\geq2$ and an empty word it is
$-c/[(n-1)(u-a)^{n-1}]$. For a one-letter word $(b)$ it is
\begin{equation}
 \mathcal P[c/(u-a)^n;(b)]
 =-\frac{cG(b;u)}{(n-1)(u-a)^{n-1}}
                  +\frac{cQ_{n-1}(u;a,b)}{n-1},
\end{equation}
where $Q_m$ is the finite elementary combination already given
above, with logarithms interpreted as the same basepoint GPLs.
Polynomial terms are fixed explicitly, for $k\geq0$, by
\begin{align}
 \mathcal P[u^k;()]&=\frac{u^{k+1}}{k+1},\\
 \mathcal P[u^k;(b)]&=\frac{u^{k+1}-b^{k+1}}{k+1}G(b;u)
  -\frac1{k+1}\sum_{j=0}^k
        \frac{b^j u^{k+1-j}}{k+1-j}.
\end{align}
Linearity completes the definition for any rational coefficient.
These operations terminate after one
step because $w$ has length at most one. Thus $\mathcal P$ denotes
a finite sum of specified GPLs and rational functions.

The complete interior primitive of the regular rational sector is
\begin{align}
 \mathcal F_{\mathrm{reg}}={}&\mathcal P[R;()]
 +\mathcal P[2K_1+K_2-K_3;(0)]
 +\mathcal P[K_2+K_3;(1)]\\
 &-\mathcal P[K_2+K_3;(s)]
 +2\mathcal P[K_3;(-s)]-2\mathcal P[K_3;(-1)].
 \label{eq:comp-real-rational-primitive}
\end{align}
Its expansion contains forty-eight GPL words of weight at most two.
This form keeps the relatively long rational coefficient table
specified by exact derivatives instead of reprinting it.

A possible internal letter in this representation is removable.
The quadratic $u^2-(s-3)u+s$ corresponds to $z=4$. The original
angular denominator $4x+z(H-x)$ is strictly positive in the physical
interior, so $z=4$ is not a singularity of the amplitude. The pole
$x=Hz/(z-4)$ merely moves through infinity at that point.
For $z<4$, the pole $Hz/(z-4)$ is negative and its partner
$H+z-Hz/(z-4)$ exceeds $H+z$. For $z>4$, the first exceeds
$H$ and the partner is smaller than $z$. Both therefore stay outside
$z<x_-\leq x\leq x_+<H$ on either side of $z=4$.
The original finite inner integral is analytic at this point.
At either root of the quadratic, $L_3=0$. Consequently the
combination $K_3L_3$ has a removable limit there, even when one root
lies inside the real $u$ interval. The GPL combination must be
formed before that limit is taken. Passing above or below the
individual artificial letter gives the same complete expression,
because its would-be discontinuity is proportional to this zero.
This argument fixes the continuation from the original real
integrand and does not add a physical threshold at $s=9$.

\paragraph{Exchange symmetry after the local subtraction.}
The full angular expression is invariant under the exchange of the
identical final photons. This exchange sends $x$ to $y=H-x+z$ and
leaves $z$ unchanged. The Taylor polynomials used to remove the Gram
zero were instead expanded at fixed $x$. These operations do not
commute beyond the order of the subtraction. In particular, direct
substitution gives
\begin{equation}
 \begin{split}
 T_{UW}(y,z)-T_{VW}(x,z)
 ={}&-\frac{s z^2}{6H(x-H)^2(x-H-z)^2}\\
 &\times[-H^2-6H+x^2-xz+6x-6z].
 \end{split}
\end{equation}
Using the coefficient $R_{VW}$ of $J_{VW}$, define the rational
exchange correction
\begin{equation}
 D(x,z)=R_{VW}(x,z)[T_{UW}(H-x+z,z)-T_{VW}(x,z)].
\end{equation}
The factor $z^2$ cancels the apparent Gram pole of $R_{VW}$.
The symmetric regulator $(xy)^\eta$, the Dalitz domain, and its
measure are all invariant under photon exchange. Changing variables
in the subtracted $VW$ integral therefore proves
\begin{equation}
 I_{VW}^{\mathrm{sub}}(\eta)
 =I_{UW}^{\mathrm{sub}}(\eta)+I_D(\eta).
\end{equation}
This identity holds first where the regulated integrals converge
and then for their meromorphic continuations. It applies separately
to the pole and finite Laurent coefficients. Consequently the real
integral is assembled from the regular rational term, the two
single-denominator logarithms, the $W$ logarithm, the $UV$ term,
twice the $UW$ term, and $D$. Omitting $D$ would change both the
infrared pole and the finite answer.

\paragraph{Endpoint values of the primitives.}
The preceding formulas specify finite combinations of GPLs rather
than unevaluated phase-space integrals. Their endpoint subtraction
can also be stated algebraically. Write $b=s-u$, let $f(u)$ denote
one complete sector integrand, and let $\mathcal F'(u)=f(u)$.
After combining all terms of that sector its expansion is
\begin{equation}
 f(s-b)=\frac{A}{b}
   +\sum_{n=0}^{\infty}(c_n+d_n\ln b)b^n.
 \label{eq:comp-real-endpoint-series}
\end{equation}
There are no higher negative powers or $\ln b/b$ terms in these
three sectors. This follows by expanding the displayed rational
coefficients and the logarithms before taking the endpoint. The
coefficient $A$ vanishes for the $W$ sector. All series coefficients
are obtained by polynomial division and the Taylor series of
$\ln(1-b/a)$.

Let $\mathcal F_0$ be the finite limit of the complete primitive at
$u=\sqrt{s}$. It includes constant contributions from products of a
vanishing GPL and a rational coefficient with a pole. They are
obtained by multiplying their local series before setting
$u=\sqrt{s}$.

The lower endpoint constants simplify considerably. Direct local
expansion of the finite GPL combinations gives
\begin{align}
 \mathcal F_{0,W}&=\mathcal F_{0,\ln s}=0,\\
 \mathcal F_{0,\mathrm{reg}}
 &=\frac{2s[-s^5-5s^4+18s^3+122s^2+151s-45]}
              {3(s^2+2s-3)}.
\end{align}
For the $W$ primitive, the coefficients that have a simple pole at
the basepoint multiply weight-two words. Those words vanish
quadratically, so they do not contribute a constant. For the
regular rational primitive the nonempty-word coefficients are
regular there. Its lower endpoint constant is therefore the
rational empty-word coefficient evaluated at $u=\sqrt{s}$.

Since
\begin{equation}
 z=\frac{H}{s}b+O(b^2),
\end{equation}
the finite part defined by a cut on $z$ is exactly
\begin{align}
 F_z={}&\mathcal F(s-\delta)-\mathcal F_0
       +A\ln\delta+A\ln(H/s)\\
 &+\sum_{n=0}^{\infty}\delta^{n+1}
  \left[\frac{c_n}{n+1}
       +d_n\left(\frac{\ln\delta}{n+1}
                         -\frac{1}{(n+1)^2}\right)\right].
 \label{eq:comp-real-gpl-endpoint-finite}
\end{align}
Here $\delta>0$ lies inside the convergence disk of the endpoint
series. The series is simply the analytic tail of the primitive.
It is not a numerical phase-space integration. Equivalently,
$F_z=\lim_{\epsilon\to0}[\int_{z>\epsilon}f\,du+A\ln\epsilon]$.
The subsequent conversion to $(xy)^\eta$ must still be included.

The artificial interior GPL letters can be bypassed during series
continuation without choosing a new physical branch. For the
regular rational sector the original finite $x$ integral is analytic
at $z=4$: its denominators have no zero there for interior $x$.
Differentiation of the specified primitive gives this same analytic
function, including at the removable limit. Its integral around a
small closed contour surrounding an artificial letter is therefore
zero. Continuations above and below it agree when the entire GPL
combination is formed. The starting value is fixed by
$\mathcal F_0$ on the real interval. This proves path independence
of the complete expression without dropping an imaginary part of
an individual GPL.

For $1<s<9$ the additional roots are complex and do not meet the
integration interval. For $s>9$ the only root that can meet it is
\begin{equation}
 q_+=\frac{s-3+\sqrt{(s-1)(s-9)}}{2}.
\end{equation}
The root at $s=9$ coincides with the lower endpoint. The original
$x$ integral remains regular there, and the coincident GPL words
are combined before their basepoint series is evaluated. Thus the
same analytic continuation covers both sides of $s=9$.

Exact symbolic differentiation verifies the displayed GPL
primitives for all three sectors.
